% SPDX-License-Identifier: CC-BY-SA-4.0
% Copyright (c) 2025 Luis Alonso
%
% This file is part of luis_alonso/Notas-Japones.
%
% Licensed under Creative Commons Attribution-ShareAlike 4.0 International (CC BY-SA 4.0).
% You may share and adapt this material for any purpose,
% as long as you provide attribution and distribute adaptations under the same license.
%
% License: https://creativecommons.org/licenses/by-sa/4.0/
% Source:  https://git.lalonso.com/luis_alonso/Notas-Japones
%
% Note: Third-party content (if any) is excluded from this license and is marked where used.

\documentclass[a4paper,lualatex,ja=standard,base=10pt]{bxjsarticle}

\title{\ruby{漢|字}{かん|じ} LV2}
\author{アロンゾ　ルイス}

\setmainjfont{Noto Serif CJK JP}
\usepackage{luatexja-fontspec}
\usepackage{kanjicard} % Custom Kanji card
\usepackage{fancyhdr}

% --- Page style ---
\pagestyle{fancy}
\fancyhf{} % clear header/footer

% Footer: name + copyright on left, page number on right
\fancyfoot[L]{Luis Alonso\ \copyright\ \the\year}
\fancyfoot[R]{\thepage}

% Optional: remove header rule and keep a footer rule
\renewcommand{\headrulewidth}{0pt}
\renewcommand{\footrulewidth}{0.4pt}

% Increase space reserved for header area
\setlength{\headheight}{0pt}
\setlength{\headsep}{6pt} % distance from header to text

\fancypagestyle{plain}{%
  \fancyhf{}
  \fancyfoot[L]{Luis Alonso\ \copyright\ \the\year}
  \fancyfoot[R]{\thepage}
  \renewcommand{\headrulewidth}{0pt}
  \renewcommand{\footrulewidth}{0.4pt}
}

\begin{document}
\setcounter{page}{1} % start page number here

\raggedbottom
\maketitle{}
\thispagestyle{fancy}

\begin{kanjicard}{広}{ひろ・い}{コウ}
  \KExample{広い}{ひろ・い}{Amplio}
  \KExample{広場}{ひろ・ば}{Espacio abierto}
\end{kanjicard}

\begin{kanjicard}{狭}{せま・い}{キョウ/コウ}
  \KExample{狭い}{せま・い}{Angosto}
  \KExample{狭量}{きょう・りょう}{Cerrado de mente, intolerante}
\end{kanjicard}

\begin{kanjicard}{明}{あか・るい}{メイ}
  \KExample{明るい}{あか・るい}{Iluminado, persona feliz}
\end{kanjicard}

\begin{kanjicard}{暗}{くら・い}{アン}
  \KExample{暗い}{くら・い}{Oscuro}
  \KExample{暗号}{あん・ごう}{Contraseña}
\end{kanjicard}

\begin{kanjicard}{大}{おお・きい}{ダイ/タイ}
  \KExample{大きい}{おお・きい}{Grande}
  \KExample{大切}{たい・せつ}{Importante}
  \KExample{大好き}{だい・すき}{Gustar mucho}
\end{kanjicard}

\begin{kanjicard}{小}{ちい・さい}{ショウ}
  \KExample{小さい}{ちい・さい}{Pequeño}
  \KExample{小説}{しょう・せつ}{Novela}
\end{kanjicard}

\begin{kanjicard}{新}{あたら・しい}{シン}
  \KExample{新しい}{あたら・しい}{Nuevo}
  \KExample{新聞}{しん・ぶん}{Periodico}
\end{kanjicard}

\begin{kanjicard}{古}{ふる・い}{コ}
  \KExample{古い}{ふる・い}{Viejo}
  \KExample{中古}{ちゅう・こ}{Segunda mano}
\end{kanjicard}

\end{document}
